\clearpage
\fancyhead[L]{\small CONNAISSANCE DES TEMS.}
\SourceLeaf{620}
\begin{center}\small Göttingische gelehrte Anzeigen. No. 191. Pages 1905..1910. 29 November 1819.\\[3pt]\rule{.72\linewidth}{.3pt}\end{center}

\noindent\emph{Connaissance des tems, ou des mouvemens célestes, à l'usage des astronomes et des navigateurs, pour l'an {\upshape 1821}. Publiée par le bureau des longitudes.} Paris 1818. Published by the widow Courcier. 336 pages in large octavo.

\begingroup\brokenpenalty=10000\relax
The appendix to the present volume of the Connaissance des tems stands out in particular for several substantial papers by France’s foremost geometers. First come the results of a theoretical investigation of the libration of the Moon by \textsc{Poisson}. More recent observations of the lunar spot Manilius by \textsc{Bouvard} are also mentioned here, and are said to give the inclination of the lunar equator as $1^\circ\,27'\,40''$ (\textsc{Tobias Mayer} had found $1^\circ\,29'$ in 1749 from observations of the same spot). It would undoubtedly be highly desirable to make more observations of this kind. --- New method for determining cometary orbits, by \textsc{Lagrange}. This paper had already appeared in German in the Berliner astronomisches Jahrbuch for 1783. The method requires six complete observations, which must be close together in pairs. If the orientation of the orbital plane is known, the heliocentric position can be derived from the geocentric position, and two heliocentric positions close to each other are known to give a very simple determination of half the parameter. It is therefore clear that each pair of nearby geocentric positions gives an equation of condition relating the inclination of the orbit, the longitude of the ascending node, and half the parameter. Thus the six complete observations give three such equations of condition, and eliminating two of the unknown quantities finally leads to an equation of the seventh degree. Given the present state of this branch of astronomical computation, there will be no occasion to use this method. --- On the rotation of the Earth, by \textsc{Laplace}. This paper contains several interesting developments concerning the free axis of rotation of the body of the Earth, insofar as the body is covered by a liquid that is in equilibrium at its own surface, and concerning the change in the position of that axis caused by the attraction of the Moon and the Sun. --- The paper that follows, by \textsc{Poisson}, on the precession of the equinoxes likewise treats this last subject in a similar way, with the elegance of exposition one has come to expect from this excellent geometer. --- On the influence of the great equation of Jupiter and Saturn on the motion of the bodies of the solar system, by \textsc{Laplace}. This influence is found to be appreciable only in the satellites of those planets themselves. There it appears as an inequality of much the same form as the great equations themselves. For Jupiter’s fourth satellite, \textsc{Laplace} finds a maximum value of 14 seconds. --- Geographical position of a large number of points in the Mediterranean, determined in the years 1816 and 1817 by frigate captain \textsc{Gauttier}. --- On the law of gravity when the terrestrial spheroid is assumed to be homogeneous and of the same density as the sea, by \textsc{Laplace}. It is remarkable that under this assumption, which admittedly does not hold in nature, gravity at the surface of the Earth, whether at sea, on an island, or on the mainland, becomes independent of the shape of the land rising above the sea. --- According to a later remark by the author, this result is independent of the assumption that the terrestrial spheroid has the \emph{same} density as the sea, and presupposes \emph{only} that the density of the spheroid is uniform. And since the formula obtained for gravity gives an appreciably smaller increase from the equator to the pole than the observed behavior of pendulum oscillations indicates, this proves that the density of the Earth, \SourceLeaf{621}from the surface toward the center, increases. Indeed, this increase in density must extend to a considerable depth. --- On the positions of the fixed stars for which lunar distances are given in the ephemerides, by \textsc{Bouvard}. From the \textsc{Bessel}-\textsc{Bradley} catalogue for 1755 and the \textsc{Piazzi} catalogue for 1800 (with the right ascensions in the latter, however, slightly altered, as stated, on the basis of observations made in Paris), the positions are carried forward to the year 1830. This reduction, however, does not seem to have been calculated accurately throughout. --- Investigations on four earlier comets, by \textsc{Burckhardt}. For the second comet of 1766, Mr. \textsc{Burckhardt} takes into account the observations made by \textsc{Lanux} on the island of Bourbon and finds an elliptical orbit with a period of 5 years. The motions of the comets of 1774, 1723, 1729 are better represented by hyperbolic orbits than by parabolic ones. --- By the same astronomer, a table for determining the time in parabolic motion from two distances from the Sun and the angle between them. Also, critical remarks on various stars observed only once by \textsc{Bradley}. --- On the English expedition to the North Pole. --- Extract from the account of Captain \textsc{Tuckey}'s expedition to explore the river Zaire. --- On the shape of the Earth and the law of gravity at its surface, by \textsc{Laplace}. The text printed here gives an overview of the main results of a paper that the author read before the Royal Academy. We have a complete printed copy of the paper before us. It treats the problem more fully and more naturally than it has ever been treated before. The novelty and fruitfulness of the methods it uses, together with the importance of the results for geology, make it one of this great geometer's finest works. --- Extract from a letter from Mr. W. \textsc{Lambton} to Mr. \textsc{Delambre} concerning the East Indian arc measurement, which at the time (in September 1817) already covered a meridian arc of nearly 10 degrees and was still being extended. Mr. \textsc{Lambton} gives the latitude of the southern endpoint Runnä as $8^\circ\,9'\,38''4$, that of the northern endpoint Darumgedda as $18^\circ\,3'\,23''26$, and the length of the arc as 598610 fathoms at $62^\circ$ Fahrenheit. By comparison with the French arc measurement, he finds a flattening of $\frac{1}{309.5}$, from the English one $\frac{1}{313.54}$, from the Swedish one $\frac{1}{307.15}$; by contrast, when he compares the whole arc with the whole arc from Dünkirchen to Montjouy, he finds, using a formula by \textsc{Playfair}, $\frac{1}{355}$ (this marked discrepancy arises from an incorrect conversion of the French arc from toises to fathoms). --- Extract giving some determinations of positions in Sweden, from the papers of the Stockholm Academy. --- On the comet of 1818, by \textsc{Nicollet}. --- Extract from \textsc{Lemonnier}'s observing journal, which contains no fewer than 12 observations of Uranus, the first from 14 October 1750, the last from 18 December 1771. In 1769, over 8 days, from 15 to 23 January, this astronomer had unknowingly observed the planet 6 times, and one can hardly understand how he himself could have missed the discovery. Mr. \textsc{Bouvard}, who brings these observations to our attention, is working on new tables for Uranus. Still, given the condition of the instrument and the clock and the observer's lack of care, the observations themselves are of little value. --- On refraction at low altitudes, by \textsc{Groombridge}. This is a supplement to an earlier paper by this astronomer, in which he had given only the mean deviation from his tables for each individual star. Here the extreme differences are reported as well. This can help to some extent in judging how closely the observations agree with one another. Yet on the important question on which astronomers have so far been divided, whether it is appropriate to use the external thermometer or the internal thermometer or a mean of the two in the reduction, it would be desirable to publish all the individual observations. The observations of four solstices from 1816---\SourceLeaf{622}1818 are also included. Here the mean of the two determinations of the obliquity of the ecliptic at the two solar solstices does not differ appreciably from the mean for the two winter solstices, and the mean of them all gives, for the time of the autumnal equinox of 1817, a mean obliquity of $23^\circ\,27'\,48''29$.
\par\clearpage\endgroup
